\paragraph{Character determinants and the factors at the auxiliary primes}
Partition \(I\) into the indices of the old inertial elements \(\delta_v\),
additional generators \(I_0\), and the detection set \(s\) containing
the test points. Write \(\eta_a\colon\mathcal A\to\mathcal S\) for
evaluation on the character \(\chi_a\). Apply \(\eta_a\) to the
entries of a row, forget the coordinates of old inertia, and then
apply an \(\mathcal S\)-linear map. Denote the resulting map by \(F\),
and its restriction to constant rows by
\(\Phi:\mathcal O^I\to\operatorname{im}(F)\).
For a submodule \(U\) containing \(F(N_\mathrm{row})\), set
\[
 Z(U)=\{\Psi\circ z:\ \Psi:\mathcal O^I\to\mathcal O/(p_\ast)
   \text{ is }\mathcal O\text{-linear and vanishes on }\Phi^{-1}U\}.
\]
The same three kinds of row as for the main block show that
\(c=\Psi\circ z\) is a cocycle for \(\chi_a\) satisfying the linear
relations among the families \(N_x(g)-1\).
The row \(z(\delta_v)-\mathbf e_v\) gives \(c(\delta_v)=0\), since
the old-inertia coordinates are forgotten.
The image of \((g-1)u_\gamma\) under \(\eta_a\) is
\[
 (\chi_a(g)-1)\bigl(c_\gamma z(\gamma)+z(a_{\mathrm c})\bigr).
\]
Thus \(c(\gamma)=0\) implies
\((\chi_a(g)-1)c(a_{\mathrm c})=0\) for all \(g\).
These are the conditions on the character cocycles to which the
bounds at the ordered auxiliary primes apply.
\leanlinks{\href{\upstream/grh/src/New/Congruences/FirstOrder/RibetMatrices.lean\#L564}{Character cocycles from the annihilators of the rows}}

First let \(F\) read only the \(I_0\) coordinates. Forgetting the
coordinates of \(g\in s\) forces \(c(g)=0\) for every \(c\in Z(U)\).
The set \(s\) consists of the test points, \(a_{\mathrm c}\), and the
new inertial elements and detects the character cocycles. Hence \(c=0\).
Since \(p_\ast\mathcal O^I\subset\Phi^{-1}U\), this conclusion can be
read by duality. The finitely generated \(\mathcal O\)-module
\(\mathcal O^I/\Phi^{-1}U\) is killed by \(p_\ast\). By Smith normal form,
linear forms to \(\mathcal O/(p_\ast)\) detect every cyclic summand.
Thus \(\Phi(z(g_i))\) is in \(U=F(N_\mathrm{row})\).
The tied row \(z(g_i)-\mathbf e_i\) then gives a tied row reading
as any specified unit vector on the \(I_0\) coordinates under \(\eta_a\).

The same unit rows can be used for all characters. Suppose \(b\)
gives the \(i\)-th unit row on the \(I_0\) coordinates for some
characters. For the next character choose rows \(m_j\) giving the
\(j\)-th unit rows for \(j\in I_0\), and replace \(b\) by
\(b+\sum_{j\in I_0}(\delta_{ij}-b_j)m_j\). On the new character this
evaluates to the \(i\)-th unit vector on \(I_0\). At each previously
treated character the coefficients of the added rows evaluate to
zero. Continue through all characters to obtain the common row \(\nu_i\).
\leanlinks{\href{\upstream/grh/src/New/Congruences/FirstOrder/FittingBlock.lean\#L326}{Common unit rows for the characters}
  \leanlinksep \href{\upstream/grh/src/New/Congruences/FirstOrder/FittingBlock.lean\#L700}{Duality producing unit rows from detection}}

For a fixed character, clear the \(I_0\) coordinates using the \(\nu_i\).
Write \(F_a\) for the resulting map to the remaining \(s\) coordinates,
and \(K_a=F_a(N_\mathrm{row})\subset\mathcal S^s\). Still
\(p_\ast\Phi(\mathcal O^I)\subset K_a\), and if \(y(g)=\Phi(z(g))\),
then \(y(g)\equiv\mathbf e_g\bmod K_a\) for \(g\in s\).
At each new prime \(H_j(\chi_a(\phi))\equiv1\bmod p_\ast\).
Adding an integral multiple of the tied row \(p_\ast X_j z(\phi)\)
to the original relation gives
\[
 y(\tau_j)\equiv X_j y(\phi_j),\qquad
 y(\tau_j')\equiv X_j y(\phi_j')\pmod{K_a},
\]
where primes denote inverse conjugates.

Set \(J_a=\{j\in J:\operatorname{type}(j)=a\}\) and
\(X_{<a}=\prod_{\operatorname{type}(j)<a}X_j\).
Let \(\mathfrak n\) be the ideal generated by the monomials that
do not divide \(X_J\), and put \(\mathfrak q_a=(\mathfrak n:X_{<a})\).
For \(\operatorname{type}(i)<a\), one has \(X_i\in\mathfrak q_a\),
since \(X_iX_{<a}\) is divisible by \(X_i^2\).
Start with \(U=K_a+\mathfrak q_a\mathcal S^s\).
For a submodule \(V\subseteq\mathcal S^s\), write
\(\operatorname{Fitt}(V)\) for the ideal generated by determinants
of square matrices with rows in \(V\).
This is the zeroth Fitting ideal of \(\mathcal S^s/V\).

The following comparison controls changes in this ideal.
Fix a submodule \(Z_*\) of the module of functions on \(G_E\).
We wish to include \(\Phi(u)\) in \(V\) whenever every linear form
\(\Psi\) with \(\Psi\circ z\in Z_*\) vanishes on \(u\).
If such a \(\Phi(u)\) lies outside \(V\), choose \(k\ge0\) with
\(q=\Phi(\varpi^k u)\notin V\) and \(\varpi q\in V\), where
\(\varpi\) is a uniformizer. This is possible because \(V\)
contains the image of \(p_\ast\mathcal O^I\).
For \(V'=V+\mathcal S q\), expansion by multilinearity in the rows
gives \(\varpi\operatorname{Fitt}(V')\subseteq\operatorname{Fitt}(V)\).
Terms containing the row \(q\) twice vanish, and a term containing it
once is handled by \(\varpi q\in V\). Moreover, the image of \(Z(V')\)
in the quotient by \(Z_*\) is strictly smaller than that of \(Z(V)\).
If the images were equal, any \(\Psi\) annihilating the rows of \(V\)
would differ from some \(\Psi'\) annihilating those of \(V'\) by a
form with \((\Psi-\Psi')\circ z\in Z_*\).
This implies \(\Psi(u)=\Psi'(u)\), so \(\Psi(\varpi^k u)=0\).
By the finite-quotient duality just used this implies \(q\in V\).
Each adjoining step is therefore paid for by a strict length
decrease. The valuation of the multiplier returning the enlarged
Fitting ideal to \(\operatorname{Fitt}(V)\) is at most the length
of \(Z(V)/(Z(V)\cap Z_*)\).
\leanlinks{\href{\upstream/grh/src/Lib/General/Algebra/FittingIdealsAndLengths.lean\#L671}{Controlling the change of the Fitting ideal by length}}

Apply the comparison in three stages. First require the cocycles
to vanish on the set \(A_a\subset G_E\) given by the bound at the
ordered primes. Take \(Z_*\) to be the functions vanishing there.
Vanishing then holds for the enlarged row module \(U_1\).
To return its Fitting ideal to that of \(U\), the multiplier
\(k_1\) has valuation at most the length of \(Z(U)\) modulo the
submodule vanishing on \(A_a\).

Next, in order of \(j\in J_a\), adjoin \(v_j=y(\phi_j)-y(\phi_j')\).
Write \(V\) for the row module including the previously adjoined
rows. For any \(i<j\) of the same type, \(v_i\in V\). At an earlier
type, \(X_i\in\mathfrak q_a\). Thus \(X_i v_i\in V\) in either case.
By the ties, the linear forms give equal values at the earlier
\(\tau_i,\tau_i'\), and at the earlier \(\phi_i,\phi_i'\) of the
same type. The congruence at the ordered primes gives
\(C(c(\tau_j)-c(\tau_j'))=0\). Finite-quotient duality now gives
\(C(y(\tau_j)-y(\tau_j'))\in V\), or \(CX_jv_j\in V\) by the ties.
Consequently \(CX_j\) multiplies the Fitting ideal after adjoining
\(v_j\) into the former one. For the row module \(T\) after all
these adjunctions, multiplication by \(k_1 C^{\#J_a}\prod_{j\in J_a}X_j\)
takes \(\operatorname{Fitt}(T)\) into \(\operatorname{Fitt}(U)\).

Finally set \(T_0=T+(X_j:j\in J)\mathcal S^s\).
In \(Z(T_0)\) the cocycles vanish on all the new inertia and take
equal values at \(\phi_j,\phi_j'\) for \(j\in J_a\).
Use the length comparison with \(Z_*\) the functions killed by
\(C\). The enlarged row module contains each \(C y(g)\), since
every \(C\)-torsion function evaluates to zero on \(Cz(g)\).
Using \(y(g)\equiv\mathbf e_g\), it also contains \(C\mathbf e_g\).
The determinant \(C^{\#s}\) of \(CI\) multiplied by the multiplier
\(k_2\) of the comparison belongs to \(\operatorname{Fitt}(T_0)\).
The rows adjoined to \(T\) to form \(T_0\) have zero constant
term. Therefore there is \(d\in\operatorname{Fitt}(T)\) with
constant term \(k_2 C^{\#s}\).

The length bound at the ordered primes gives \(b_a\) for the sum
of the two lengths paid for by \(k_1,k_2\). Here \(b_a\) is the
sum of the crystalline colength and the length of the odd classes.
Set \(\Gamma_a=k_1 C^{\#J_a}d\). We have
\[
 \Gamma_a\prod_{j\in J_a}X_j\in\operatorname{Fitt}(U),\qquad
 v_\ell(\Gamma_a(0))\le b_a+(\#J_a+\#s)v_\ell(C).
\]
Since \(U=K_a+\mathfrak q_a\mathcal S^s\), expansion by rows gives
\(\operatorname{Fitt}(U)\subseteq\operatorname{Fitt}(K_a)+\mathfrak q_a\).
Thus \(d_a+\xi_a=\Gamma_a\prod_{j\in J_a}X_j\) for
\(d_a\in\operatorname{Fitt}(K_a)\) and \(\xi_a\in\mathfrak q_a\).
\leanlinks{\href{\upstream/grh/src/New/Congruences/FirstOrder/FittingBlock.lean\#L124}{The three stages of the character Fitting estimate}
  \leanlinksep \href{\upstream/grh/src/New/Congruences/FirstOrder/RibetMatrices.lean\#L706}{Application to the images of the rows}}

Write \(d_a=\sum_i w_{a,i}\det B_{a,i}\).
Each row of \(B_{a,i}\) lifts to a tied row. Clear its \(I_0\)
coordinates with the \(\nu_i\), so under \(\eta_a\) the lift has
zero \(I_0\) part and the designated \(s\) part.
Choose one matrix on each character, with choice indexed by \(\mathbf i\).
There are elements \(u_a\in\mathcal A\) with \(\eta_{a'}(u_a)=C\)
if \(a'=a\) and \(0\) otherwise. Indeed, multiply
\(\prod_{a'\ne a}(a_{\mathrm c}-\chi_{a'}(a_{\mathrm c}))\) by the
quotient of \(C\) by its value at \(\eta_a\).
For each \(g\in s\), multiply the lift of its row for \(\eta_a\)
by \(u_a\) and sum over \(a\). Use the rows \(\nu_i\) for \(I_0\).
For an old inertial element \(\delta_v\), use the row consisting
of \(o_v(\delta_v)\) in its own coordinate and zero elsewhere.
The last row is tied because \((\delta_v-1)o_v(\delta_v)\) is
killed by every representation. The resulting square matrix
\(B_{\mathbf i}\) is block triangular under \(\eta_a\), with
diagonal block determinants \(\prod_v o_v(1),\,1\), and
\(C^{\#s}\det B_{a,i(a)}\). Setting \(c_0=C^{\#s}\prod_v o_v(1)\), we get
\[
 \det\eta_a(B_{\mathbf i})=c_0\det B_{a,i(a)}.
\]
\leanlinks{\href{\upstream/grh/src/New/Congruences/FirstOrder/RibetMatrices.lean\#L933}{A single tied matrix for all the characters}}
